\section{Localization and a fixed lifted diagram}
\label{sec:localization}

Let $k$ be a field, let $R$ be a finitely presented unital $k$-algebra,
let $e\in R$ be a central idempotent, and let
$\alpha:R\to R$ be a unital endomorphism.  As in the preceding section,
\[
  H(Y)={}_{\alpha}(eY)
\]
for a left $R$-module $Y$: the underlying vector space is $eY$, and
$r\in R$ acts there by $\alpha(r)$.  Centrality of $e$ ensures that this
action preserves $eY$.  The construction in this section uses these
fixed data and applies to every finite-dimensional $Y$.

We first encode the presentation of $R$ in a finite-dimensional directed
algebra.  This is a three-vertex version of the homogenized quiver
construction in \cite[Section~1, Theorem~1.1 and Proposition~1.2]{NRS2004}.
We then work directly with a Verdier
quotient and its evaluation functors.

\subsection{A quadratic presentation and its directed encoding}

Choose a presentation
\[
  R=k\langle x_1,\ldots,x_d\rangle/(p_1,\ldots,p_t)
\]
in which each $p_j$ has degree at most two, allowing constant and linear
terms.  Such a presentation exists: in any finite presentation, introduce
finitely many auxiliary generators for the successive partial products
in its finitely many monomials, together with their binary-product
definitions.  The resulting relations have degree at most two, and
eliminating the auxiliary generators recovers the original presentation.

Consider the quiver with vertices $0,1,2$ and arrow sets
\[
  \mathcal A_{01}=\{s,a_1,\ldots,a_d\},\qquad
  \mathcal A_{12}=\{s',a'_1,\ldots,a'_d\},\qquad
  \mathcal A=\mathcal A_{01}\sqcup\mathcal A_{12},
\]
where the subscripts specify the source and target.  Products of arrows
are written in the order of composition, from right to left.  Impose
the relations
\begin{equation}\label{eq:quadratic-encoding}
  a'_h s=s'a_h\quad(1\le h\le d),
  \qquad \widetilde p_j=0\quad(1\le j\le t),
\end{equation}
where homogenization replaces
\[
  x_hx_l\longmapsto a'_h a_l,\qquad
  x_h\longmapsto s'a_h,\qquad
  1\longmapsto s's.
\]
Let $B$ be the resulting path algebra with relations, and let
$\mathbf e_0,\mathbf e_1,\mathbf e_2$ be its vertex idempotents.
Write $\mathcal R$ for a fixed vector-space basis of the span of all
the relations in \eqref{eq:quadratic-encoding}, viewed in the free
space of paths of length two.  Thus each $\rho\in\mathcal R$ has a
unique expression
\[
  \rho=\sum_{\substack{a\in\mathcal A_{01}\\b\in\mathcal A_{12}}}
             \rho_{ba}\,ba.
\]
The basis $\mathcal R$ generates the same relation ideal.  There are
no paths of length three, so $B$ is finite-dimensional.

The following elementary bound will also be used for other directed
algebras later.

\begin{lemma}[Directed dimension bound]\label{lem:directed-bound}
Let $\Lambda$ be a finite-dimensional $k$-algebra with complete
orthogonal idempotents $\varepsilon_0,\ldots,\varepsilon_{n-1}$ such that
\[
  1=\sum_{i=0}^{n-1}\varepsilon_i,\qquad
  \varepsilon_i\Lambda\varepsilon_i=k\varepsilon_i,\qquad
  \varepsilon_j\Lambda\varepsilon_i=0\quad(j<i).
\]
Then $\Lambda$ has global dimension at most $n-1$ on both sides.
In particular, the algebra $B$ above has global dimension at most two
on both sides.
\end{lemma}

\begin{proof}
Say that a left module $M$ is supported at indices at least $p$ if
$\varepsilon_jM=0$ for $j<p$.  There is a projective surjection
\[
  \bigoplus_{i=p}^{n-1}
      \Lambda\varepsilon_i\otimes_k\varepsilon_iM
      \longrightarrow M,
  \qquad c\otimes v\longmapsto cv.
\]
The source is projective even when $M$ is not finitely generated,
since each summand is a direct sum of copies of
$\Lambda\varepsilon_i$.  Its components below $p$ vanish, and its
component at $p$ maps isomorphically onto $\varepsilon_pM$:
only the summand $i=p$ contributes there.  Its kernel is therefore
supported at indices at least $p+1$.

Starting with $p=0$ and repeating this construction, the
$(n-1)$st kernel is supported only at the last vertex.  Every such
module is projective, since
$\Lambda\varepsilon_{n-1}=k\varepsilon_{n-1}$.
This gives the left dimension bound, including the case $n=1$.
The same proof for $\Lambda^{\mathrm{op}}$, with the reverse ordering
of the idempotents, gives the right bound.  The path directions and
the homogeneous quadratic relations of $B$ give precisely the
displayed hypotheses with $n=3$.
\end{proof}

For a finite-dimensional left $R$-module $Y$, define $M(Y)$ to be the
left $B$-module with $Y$ at every vertex, with $s,s'$ acting as the
identity, and with $a_h,a'_h$ acting as $x_h$.  The relations of $B$
hold because they evaluate to the defining relations of $R$.

\subsection{The quotient action and evaluation}

Let
\[
  K=\K^b(\proj\text{-}B),\qquad E_i=\mathbf e_iB,
\]
so $K$ is the homotopy category of bounded complexes of finitely
generated projective \emph{right} $B$-modules.  With our path convention,
an arrow $a:i\to j$ belongs to $\mathbf e_jB\mathbf e_i$.
Left multiplication by $a$ is consequently a right-module map
$E_i\to E_j$.  Composing these maps agrees with multiplying paths
in the stated order.

Let $\mathcal S$ be the thick subcategory of $K$ generated by the
cones of $s:E_0\to E_1$ and $s':E_1\to E_2$, and form
\[
  q:K\longrightarrow Q=K/\mathcal S.
\]
We use the same notation for objects and arrows after applying $q$.
The maps $s$ and $s'$ are isomorphisms in $Q$.  Put $L=E_0$ and
\[
  \overline x_h=s^{-1}a_h\in\End_Q(L).
\]
The relation $a'_h s=s'a_h$ gives $a'_h=s'a_hs^{-1}$, whence
\[
  (s's)^{-1}a'_ha_l
     =s^{-1}a_hs^{-1}a_l
     =\overline x_h\overline x_l.
\]
Likewise $(s's)^{-1}s'a_h=\overline x_h$ and
$(s's)^{-1}s's=\id_L$.  The homogenized relations therefore give
a unital homomorphism
\begin{equation}\label{eq:quotient-action}
  \theta:R\longrightarrow\End_Q(L),
  \qquad x_h\longmapsto\overline x_h.
\end{equation}

For each finite-dimensional $Y$, tensoring defines an exact functor
\[
  \operatorname{ev}_Y:K\longrightarrow\D^b(k),
  \qquad T\longmapsto T\otimes_B M(Y).
\]
Here $\D^b(k)$ denotes the bounded derived category of finite-dimensional
$k$-vector spaces; termwise projectivity of $T$ computes the derived
tensor product.  Under the natural identifications
\[
  E_i\otimes_B M(Y)\cong\mathbf e_iM(Y)=Y,
\]
an arrow map evaluates to its action on $M(Y)$.  Thus
$\operatorname{ev}_Y$ takes $s,s'$ to identities and their cones to zero.
The kernel of an exact functor is thick, so it contains $\mathcal S$.
Equivalently, every map whose cone belongs to $\mathcal S$ evaluates
to an isomorphism.  The functor consequently factors through $Q$.
We retain the notation $\operatorname{ev}_Y$ for this factorization.
On $L$ the homomorphism \eqref{eq:quotient-action} evaluates to the
given $R$-action on $Y$.

\subsection{An odd double of a formal summand}

The idempotent $\theta(e)$ evaluates to the projection of $Y$ onto
$eY$, but it need not split inside $Q$. We therefore seek an object
of $Q$ whose evaluation is $eY$ plus a shifted copy. The following
lemma gives such an object for a copy shifted by three.

For an additive category $\mathcal T$, its idempotent completion
$\Kar(\mathcal T)$ has objects $(T,p)$, where $p^2=p\in\End(T)$,
and morphisms $u:(T,p)\to(T',p')$ satisfying $u=p'up$.
The embedding $T\mapsto(T,\id_T)$ is fully faithful.
If $\mathcal T$ has a shift functor, it extends
by $(T,p)[1]=(T[1],p[1])$.

We need the following consequence of the triangle axioms.  It is
compatible with the general idempotent-completion result of
Balmer and Schlichting \cite[Theorem~1.5]{BalmerSchlichting2001}; the argument
below uses only triangles already present in $\mathcal T$.

\begin{lemma}[Odd double]\label{lem:odd-double}
Let $\mathcal T$ be a triangulated category, let $L\in\mathcal T$,
and let $\epsilon\in\End_{\mathcal T}(L)$ be idempotent.
Put $U=(L,\epsilon)\in\Kar(\mathcal T)$.
Both $U\oplus U[1]$ and $U\oplus U[3]$ are isomorphic in
$\Kar(\mathcal T)$ to objects of $\mathcal T$.
\end{lemma}

\begin{proof}
First consider a triangle from $\mathcal T$,
\[
  A\xrightarrow{f}B\xrightarrow{g}T\xrightarrow{h}A[1],
\]
and suppose that formal decompositions
$A=I\oplus A_0$, $B=I\oplus B_0$ identify $f$ with the identity
on $I$ and zero elsewhere.  For any $Z\in\Kar(\mathcal T)$,
the long exact sequence $\Hom(Z,-)$ still holds for this triangle:
$Z$ is a retract of an object of $\mathcal T$, and its Hom sequence
is a direct summand of the exact sequence for that object.

The equalities $gf=0$ and $f[1]h=0$ give factorizations
\[
  g=g_0\operatorname{pr}_{B_0},\qquad
  h=\operatorname{in}_{A_0[1]}h_0.
\]
Exactness gives, for every such $Z$, a short exact sequence
\begin{equation}\label{eq:formal-cone-hom}
  0\longrightarrow\Hom(Z,B_0)
   \xrightarrow{(g_0)_*}\Hom(Z,T)
   \xrightarrow{(h_0)_*}\Hom(Z,A_0[1])
   \longrightarrow0.
\end{equation}
Indeed, a map into $B_0$ killed by $g_0$ factors through $f$ after
inclusion into $B$, and projection back to $B_0$ makes it zero.
At the other end, every map into $A_0[1]$ is killed by $f[1]$ after
inclusion into $A[1]$, so it lifts through $h$.
Taking $Z=A_0[1]$ supplies a section $\sigma$ of $h_0$.
It follows from \eqref{eq:formal-cone-hom} that
\[
  (g_0,\sigma):B_0\oplus A_0[1]\longrightarrow T
\]
induces isomorphisms on $\Hom(Z,-)$ for every $Z$, hence is an
isomorphism.

Apply this calculation to a triangle on $1-\epsilon:L\to L$.
Its identity summand is $(L,1-\epsilon)$, and its zero summand on
both sides is $U$.  Its cone is an object $S\in\mathcal T$ with
$S\cong U\oplus U[1]$ in the completion.  Under this identification,
consider the map
\[
  S[1]\cong U[1]\oplus U[2]
     \xrightarrow{\left(\begin{smallmatrix}0&0\\\id&0\end{smallmatrix}\right)}
  U\oplus U[1]\cong S.
\]
Fullness of $\mathcal T\hookrightarrow\Kar(\mathcal T)$ makes
this a morphism between the old objects $S[1]$ and $S$.
Its cone is again in $\mathcal T$.
The preceding calculation, with common identity summand $U[1]$,
source complement $U[2]$, and target complement $U$, identifies
that cone with $U\oplus U[3]$.
\end{proof}

Apply \Cref{lem:odd-double} in $Q$ to $\epsilon=\theta(e)$ and write
\[
  U=(L,\theta(e)),\qquad V=U\oplus U[3].
\]
Fix once and for all an object $V_0\in Q$ and an isomorphism
$V_0\cong V$ in $\Kar(Q)$ supplied by \Cref{lem:odd-double}.
Centrality of $e$ gives a unital action of $R$ on $U$ by
\[
  r\longmapsto\theta(e\alpha(r)e)\in\End_{\Kar(Q)}(U).
\]
Its unit is $\theta(e)=\id_U$, and its products agree with those
in $R$ because $e$ commutes with every $\alpha(r)$.
Use the same action on $U[3]$ and transport the resulting action
on $V$ to $V_0$.

There is now a $B$-diagram in $Q$ with $V_0$ at all three vertices:
$s,s'$ act as identities, and $a_h,a'_h$ act by the endomorphism
just defined for $x_h$.  Its relations hold by the unital
$R$-action.  Here a \emph{$B$-diagram} in an additive $k$-linear
category means objects at the three vertices and arrow maps satisfying
the relations of $B$ in that category.

The category $\D^b(k)$ splits idempotents: every bounded complex of
vector spaces is isomorphic to its cohomology with zero differential,
where splitting can be done degree by degree.  Thus evaluation extends
additively to $\Kar(Q)$.  It takes $U$ to $eY$, and the preceding
action of $r$ to the restriction of $\alpha(r)$ to $eY$.
Consequently the evaluated $B$-diagram has
\begin{equation}\label{eq:target-evaluation}
  eY\oplus(eY)[3]
\end{equation}
at every vertex, identities on $s,s'$, and the action of
$\alpha(x_h)$ on both summands for $a_h,a'_h$.

\subsection{Clearing denominators for the finite diagram}

The quotient diagram is now fixed and has the required evaluation on
every $Y$. We next represent it by one diagram in the homotopy category,
so that the same chain data will work under all those evaluations.

We give the lifting argument together with the fraction rules it uses.
Let
\[
  \Sigma=\{u\in\operatorname{Mor}(K):\Cone(u)\in\mathcal S\}.
\]
The Verdier quotient is the localization of $K$ at $\Sigma$.
The following triangle arguments give its right-fraction description;
see also \cite[Section 13.6]{StacksProject}.
Identities belong to $\Sigma$, and composition preserves $\Sigma$
by the octahedral axiom.  If $u:Y'\to Y$ belongs to $\Sigma$ and
$f:X\to Y$ is any map, complete $u$ to a triangle
\[
  Y'\xrightarrow{u}Y\xrightarrow{c}T\longrightarrow Y'[1],
  \qquad T\in\mathcal S.
\]
Complete $cf$ to a triangle
$X'\xrightarrow{v}X\xrightarrow{cf}T\to X'[1]$.
Then $v\in\Sigma$, and exactness of $\Hom(X',-)$ gives
$g:X'\to Y'$ with $fv=ug$.  This is the required Ore square.

For the cancellation rule, suppose that $u:Y\to Y'$ belongs to
$\Sigma$ and $ug=0$ for $g:X\to Y$.
A triangle $T[-1]\to Y\xrightarrow{u}Y'\to T$ shows that $g$
factors through a map $b:X\to T[-1]$ with $T\in\mathcal S$.
Completing $b$ to a triangle gives a denominator $v:X'\to X$
with $bv=0$, hence $gv=0$.  Applying this to differences gives
the corresponding rule for equality of morphisms.

These rules construct the localization by roofs
\[
  X\xleftarrow{u}X'\xrightarrow{f}Y,
  \qquad u\in\Sigma,
\]
representing $q(f)q(u)^{-1}$; equality is tested after a common source
refinement.  In particular, every morphism $q(X)\to q(Y)$ has this
form, and
\begin{equation}\label{eq:zero-fraction}
  q(g)=0\quad\Longrightarrow\quad
  gv=0\text{ for some }v:X'\to X\text{ in }\Sigma
\end{equation}
for $g:X\to Y$.  For the last assertion compare the roofs
$(\id_X,g)$ and $(\id_X,0)$.
Finite families can be treated with one denominator by collecting
their maps into a single map to a finite direct sum of targets.

\begin{lemma}[Finite diagram lifting]\label{lem:diagram-lift}
Every $B$-diagram in $Q$ is isomorphic to the image under $q$ of a
$B$-diagram in $K$.
\end{lemma}

\begin{proof}
Write the given objects as $W_i$ and its arrow maps as $t_a$.
Choose $C_i\in K$ and isomorphisms $\psi_i:q(C_i)\to W_i$.
First collect the maps
\[
  (\psi_2^{-1}t_b\psi_1)_{b\in\mathcal A_{12}}:
     q(C_1)\longrightarrow
     \bigoplus_{b\in\mathcal A_{12}}q(C_2).
\]
Represent this tuple by one roof, with denominator
$u_1:D_1\to C_1$ and numerator components $f_b:D_1\to C_2$.
Set $D_2=C_2$ and
\[
  \phi_1=\psi_1q(u_1),\qquad \phi_2=\psi_2.
\]
Then $\phi_2q(f_b)=t_b\phi_1$ for every $b\in\mathcal A_{12}$.

Next collect the adjusted maps
\[
  (\phi_1^{-1}t_a\psi_0)_{a\in\mathcal A_{01}}:
     q(C_0)\longrightarrow
     \bigoplus_{a\in\mathcal A_{01}}q(D_1).
\]
A single roof gives a denominator $u_0:D_0\to C_0$ and maps
$f_a:D_0\to D_1$.  Put $\phi_0=\psi_0q(u_0)$.
All the arrow squares now commute:
\begin{equation}\label{eq:lift-arrow-squares}
  \phi_jq(f_a)=t_a\phi_i\qquad(a:i\to j).
\end{equation}

For $\rho\in\mathcal R$, form the morphism in $K$
\[
  g_\rho=\sum_{b,a}\rho_{ba}f_bf_a:D_0\longrightarrow D_2.
\]
The relations for $t_a$ and \eqref{eq:lift-arrow-squares} give
$\phi_2q(g_\rho)=0$, so $q(g_\rho)=0$.
Apply \eqref{eq:zero-fraction} to the tuple $(g_\rho)_\rho$.
There is one denominator $v:D'_0\to D_0$ such that
$g_\rho v=0$ in $K$ for every $\rho$.
Replace $D_0$ by $D'_0$, every map $f_a$ with source vertex $0$ by
$f_av$, and $\phi_0$ by $\phi_0q(v)$.
The arrow squares still commute, and each relation becomes
\[
  \sum_{b,a}\rho_{ba}f_b(f_av)=g_\rho v=0
  \quad\text{in }K.
\]
There are no arrows into vertex zero and no longer paths requiring
further compatibility.  Thus the new data form a $B$-diagram in $K$,
and the three $\phi_i$ give an isomorphism from its image to the
entire prescribed diagram in $Q$.
\end{proof}

\subsection{The chain data for tensor realization}

Apply \Cref{lem:diagram-lift} to the diagram on $V_0$ constructed
above.  The resulting relations hold in the homotopy category.
Choosing representatives converts them into a finite list of
specified null-homotopies, which are the input for the next section.

\begin{proposition}[Fixed lifted diagram]\label{prop:lifted-diagram}
For the fixed data $R,e,\alpha$ and the algebra $B$ above, there are
bounded complexes $D_0,D_1,D_2$ of finitely generated projective
right $B$-modules, chain maps
\[
  f_a:D_i\longrightarrow D_j\qquad(a:i\to j),
\]
and right-$B$-linear maps $h_\rho:D_0\to D_2$ of degree $-1$
for $\rho\in\mathcal R$, such that
\begin{equation}\label{eq:relation-homotopies}
  d_{D_2}h_\rho+h_\rho d_{D_0}
      =\sum_{b,a}\rho_{ba}f_bf_a.
\end{equation}
For every finite-dimensional left $R$-module $Y$, the diagram
whose objects and arrows in $\D^b(k)$ are
\begin{equation}\label{eq:evaluated-diagram}
  D_i\otimes_B M(Y),\qquad f_a\otimes_B\id_{M(Y)}
\end{equation}
is isomorphic to the diagram having $eY\oplus(eY)[3]$ at all three
vertices, identities on $s,s'$, and the action of $\alpha(x_h)$
on both summands for $a_h,a'_h$.
All the complexes, maps, and homotopies in this statement are fixed
before $Y$ is chosen.
\end{proposition}

\begin{proof}
The presentation, the quotient, the two cones, their isomorphisms in
the completion, and the action on $V_0$ have all been chosen using
only $R,e,\alpha$.  Apply the finite lifting construction once to
this fixed diagram.  Its objects lie in $K$, so choose bounded
right-projective complex representatives $D_i$ and chain-map
representatives $f_a$.  Each relation is zero in $K$, hence its
representative is null-homotopic.  Choose one right-linear
homotopy $h_\rho$ for each of the finitely many relations.
This gives \eqref{eq:relation-homotopies}.

The isomorphisms $\phi_i$ from the lifting construction are fixed
morphisms in $Q$ and satisfy all the fixed arrow squares
\eqref{eq:lift-arrow-squares}.  Every $\operatorname{ev}_Y$ factors
through this same quotient.  Applying it to those squares and then
to the fixed identification $V_0\cong V$ gives the asserted
isomorphism of diagrams, by \eqref{eq:target-evaluation}.
Thus the same choices work for all $Y$.
\end{proof}

The construction has used the homomorphism
\eqref{eq:quotient-action} and the evaluation functors alone.
It requires neither an identification of $Q$ with a derived category
of $R$ nor a stable-flatness assertion about a universal localization.
The maps $f_a$ satisfy the relations up to the chosen homotopies
\eqref{eq:relation-homotopies}; the next section incorporates those
homotopies into an actual complex of bimodules.
